\documentclass[12pt,a4paper]{article}

\usepackage[margin=1.4cm]{geometry}
\usepackage{array}

\usepackage{enumitem}

% Global compact spacing for lists
\setlist[itemize]{
	topsep=0pt,
	partopsep=0pt,
	parsep=0pt,
	itemsep=0pt
}

\setlist[enumerate]{
	topsep=0pt,
	partopsep=0pt,
	parsep=0pt,
	itemsep=0pt
}

\usepackage{multicol}

\newcolumntype{L}[1]{>{\raggedright\arraybackslash}p{#1}}

\newcommand{\checkbox}{\fbox{\rule{0pt}{0.8ex}\rule{0.8ex}{0pt}}}

\usepackage{xcolor}
\definecolor{solutionred}{HTML}{7F1D1D}
\newcommand{\sol}[1]{{\color{solutionred}\small\strut #1\strut}}
\newcommand{\answerline}[1]{\makebox[0pt][l]{\raisebox{0.6ex}{\sol{#1}}}}
\begin{document}
	

	\begin{center}
		{\Large\bfseries C1: OPORTUNIDADES A NUESTRO ALREDEDOR}
		
		\vspace{0.1cm}
		
		Emprendimiento e Innovación Grado 10\\[-1mm]{\scriptsize\color{solutionred}Modelo pedagógico: toda evidencia y respuesta es simulada; no son datos del colegio.}
		
		\vspace{0.3cm}
		
		\begin{tabular}{L{0.13\textwidth}L{0.65\textwidth}}
			Grado: & \answerline{10}\rule{2cm}{0.4pt} \\
			Fecha: & \rule{2cm}{0.4pt} \\
			Nombre 1: & \rule{\linewidth}{0.4pt} \\
			Nombre 2: & \rule{\linewidth}{0.4pt} \\
			Nombre 3: & \rule{\linewidth}{0.4pt} \\
		\end{tabular}
	\end{center}
	
	\section*{Nuestro objetivo}
	
	La innovación comienza cuando identificamos algo que podría mejorarse. Primero, generarán ideas a partir de su experiencia en el aula. Luego, como tarea, observarán su entorno e identificarán situaciones que puedan representar una necesidad, un problema, una incomodidad o una oportunidad de mejora.
	
	En esta etapa, no diseñen una solución. El objetivo es identificar situaciones que valga la pena investigar.
	
	\section*{Actividad en clase: Generar ideas}
	
	Piensen en su experiencia en el colegio y en el aula. Consideren situaciones que:
	
	\noindent\begin{tabular}{@{}p{0.43\linewidth}@{\hspace{0.04\linewidth}}p{0.53\linewidth}@{}}
		\begin{itemize}[leftmargin=0.5cm,itemsep=0pt]
			\item Hacen que una actividad sea difícil o poco eficiente.
			\item Generan incomodidad o frustración.
			\item Desperdician tiempo, materiales o recursos.
		\end{itemize}
		&
		\begin{itemize}[leftmargin=0.5cm,itemsep=0pt]
			\item Generan dificultades para estudiantes, profesores u otras personas.
			\item Podrían organizarse o mejorarse.
			\item Podrían realizarse de una manera diferente.
		\end{itemize}
	\end{tabular}
	
	\vspace{0.3cm}
	
	\textbf{Ideas de nuestro salón de clase:}
	
	\vspace{0.1cm}
	
	1.\ \noindent\answerline{Filas para comprar alimentos durante el descanso.}\rule{0.85\linewidth}{0.4pt}
	
	\vspace{0.25cm}
	
	2.\ \noindent\answerline{Ruido que dificulta escuchar instrucciones en trabajos grupales.}\rule{0.85\linewidth}{0.4pt}
	
	\vspace{0.25cm}
	
	3.\ \noindent\answerline{Materiales compartidos difíciles de localizar al iniciar actividades.}\rule{0.85\linewidth}{0.4pt}
	
	\vspace{0.25cm}
	
	4.\ \noindent\answerline{Dudas sobre qué versión de una tarea se debe entregar.}\rule{0.85\linewidth}{0.4pt}
	
	\vspace{-5mm}
	\section*{Observen su entorno}
	
	Durante su día normal, presten atención a lo que sucede a su alrededor. Pueden observar:
	
	\begin{multicols}{2}
		\begin{itemize}[leftmargin=0.6cm,itemsep=0pt]
			\item Su colegio
			\item Su hogar
			\item Su barrio
			\item Espacios públicos
			\item Transporte
			\item Tiendas o negocios
			\item Espacios recreativos
			\item Situaciones relacionadas con la tecnología
			\item Otros entornos de la comunidad.
		\end{itemize}
	\end{multicols}
	
	Busquen situaciones en las que las personas experimenten una dificultad, necesidad, incomodidad, ineficiencia o una oportunidad que no esté siendo atendida. \vspace{6mm}
	
	\noindent\begin{tabular}{@{}p{0.32\linewidth}p{0.63\linewidth}@{}}
		No comiencen preguntándose:
		&
		\textit{``¿Qué debería construir?''}
		\\[0.15cm]
		&\\
		En cambio, pregúntense:
		&
		\textit{``¿Qué está sucediendo y qué podría mejorarse?''}
	\end{tabular}
	
	\newpage
	
	\section*{Qué deben registrar}
	
	Para cada oportunidad, registren:
	
	\begin{center}
		\small
		\begin{tabular}{|L{0.20\textwidth}|L{0.26\textwidth}|L{0.38\textwidth}|}
			\hline
			\textbf{Elemento} & \textbf{Qué escribir} & \textbf{Ejemplo} \\
			\hline
			\textbf{Lugar} & ¿Dónde? & Cafetería del colegio \\
			\hline
			\textbf{Situación} & ¿Qué está sucediendo? & Los estudiantes esperan en una fila larga. \\
			\hline
			\textbf{Necesidad / Problema / Incomodidad} & ¿Qué necesita mejorar? & Los estudiantes pasan tiempo esperando. \\
			\hline
			\textbf{¿A quién afecta?} & ¿Quién lo experimenta? & Estudiantes y personal de la cafetería. \\
			\hline
			\textbf{Evidencia} & ¿Qué observaron? & La fila se hizo muy larga. \\
			\hline
			\textbf{¿Por qué importa?} & ¿Por qué vale la pena investigarlo? & Los estudiantes tienen menos tiempo para comer. \\
			\hline
			\textbf{Posible solución} & ¿Ya tienen una? & Dejar abierta por ahora. \\
			\hline
		\end{tabular}
	\end{center}
	
	\textbf{Importante:} Primero comprendan la situación. No decidan la solución demasiado pronto.
	
	\section*{Oportunidad 1}
	
	\begin{center}
		\small
		\begin{tabular}{|L{0.22\textwidth}|L{0.61\textwidth}|}
			\hline
			\textbf{Lugar} & \rule{0pt}{0.5cm}\begin{minipage}[b]{\linewidth}\sol{Cafetería: punto de venta de alimentos durante el descanso.}\end{minipage} \\
			\hline
			\textbf{Situación} & \rule{0pt}{0.5cm}\begin{minipage}[b]{\linewidth}\sol{Varios estudiantes esperan para comprar; la fila parece extensa.}\end{minipage} \\
			\hline
			\textbf{Necesidad / Problema / Incomodidad} & \rule{0pt}{0.5cm}\begin{minipage}[b]{\linewidth}\sol{La espera podría reducir el tiempo para comer y descansar.}\end{minipage} \\
			\hline
			\textbf{¿A quién afecta?} & \rule{0pt}{0.45cm}\begin{minipage}[b]{\linewidth}\sol{Estudiantes compradores y personal que los atiende.}\end{minipage} \\
			\hline
			\textbf{Evidencia} & \rule{0pt}{0.6cm}\begin{minipage}[b]{\linewidth}\sol{Escena preliminar simulada: varios esperan mientras otros llegan.}\end{minipage} \\
			\hline
			\textbf{¿Por qué importa?} & \rule{0pt}{0.5cm}\begin{minipage}[b]{\linewidth}\sol{El descanso es limitado; falta verificar cuánto ocupa la compra.}\end{minipage} \\
			\hline
			\textbf{Posible solución} & \rule{0pt}{0.45cm}\begin{minipage}[b]{\linewidth}\sol{Abierta: primero investigar la espera y sus posibles causas.}\end{minipage} \\
			\hline
		\end{tabular}
	\end{center}
	
	\vspace{0.15cm}
	
	\section*{Oportunidad 2}
	
	\begin{center}
		\small
		\begin{tabular}{|L{0.22\textwidth}|L{0.61\textwidth}|}
			\hline
			\textbf{Lugar} & \rule{0pt}{0.5cm}\begin{minipage}[b]{\linewidth}\sol{Salón de décimo durante el trabajo en grupos.}\end{minipage} \\
			\hline
			\textbf{Situación} & \rule{0pt}{0.5cm}\begin{minipage}[b]{\linewidth}\sol{Varias conversaciones coinciden con las instrucciones.}\end{minipage} \\
			\hline
			\textbf{Necesidad / Problema / Incomodidad} & \rule{0pt}{0.5cm}\begin{minipage}[b]{\linewidth}\sol{Escuchar y comprender las indicaciones sin repetirlas.}\end{minipage} \\
			\hline
			\textbf{¿A quién afecta?} & \rule{0pt}{0.45cm}\begin{minipage}[b]{\linewidth}\sol{Estudiantes y docente.}\end{minipage} \\
			\hline
			\textbf{Evidencia} & \rule{0pt}{0.6cm}\begin{minipage}[b]{\linewidth}\sol{Escena simulada: dos grupos piden repetir una indicación.}\end{minipage} \\
			\hline
			\textbf{¿Por qué importa?} & \rule{0pt}{0.5cm}\begin{minipage}[b]{\linewidth}\sol{Puede dificultar el inicio de la actividad.}\end{minipage} \\
			\hline
			\textbf{Posible solución} & \rule{0pt}{0.45cm}\begin{minipage}[b]{\linewidth}\sol{Abierta; solo se registra como oportunidad inicial.}\end{minipage} \\
			\hline
		\end{tabular}
	\end{center}
	
	\vspace{0.15cm}
	
	\section*{Oportunidad 3}
	
	\begin{center}
		\small
		\begin{tabular}{|L{0.22\textwidth}|L{0.61\textwidth}|}
			\hline
			\textbf{Lugar} & \rule{0pt}{0.5cm}\begin{minipage}[b]{\linewidth}\sol{Zona de materiales compartidos del salón.}\end{minipage} \\
			\hline
			\textbf{Situación} & \rule{0pt}{0.5cm}\begin{minipage}[b]{\linewidth}\sol{Se busca una regla antes de iniciar una actividad.}\end{minipage} \\
			\hline
			\textbf{Necesidad / Problema / Incomodidad} & \rule{0pt}{0.5cm}\begin{minipage}[b]{\linewidth}\sol{Localizar los materiales disponibles con claridad.}\end{minipage} \\
			\hline
			\textbf{¿A quién afecta?} & \rule{0pt}{0.45cm}\begin{minipage}[b]{\linewidth}\sol{Estudiantes que comparten útiles.}\end{minipage} \\
			\hline
			\textbf{Evidencia} & \rule{0pt}{0.6cm}\begin{minipage}[b]{\linewidth}\sol{Escena simulada: un grupo revisa varias cajas.}\end{minipage} \\
			\hline
			\textbf{¿Por qué importa?} & \rule{0pt}{0.5cm}\begin{minipage}[b]{\linewidth}\sol{La búsqueda puede restar tiempo de trabajo.}\end{minipage} \\
			\hline
			\textbf{Posible solución} & \rule{0pt}{0.45cm}\begin{minipage}[b]{\linewidth}\sol{Abierta; no se desarrolla un proyecto adicional.}\end{minipage} \\
			\hline
		\end{tabular}
	\end{center}
	
	\section*{Oportunidad 4}
	
	\begin{center}
		\small
		\begin{tabular}{|L{0.22\textwidth}|L{0.61\textwidth}|}
			\hline
			\textbf{Lugar} & \rule{0pt}{0.5cm}\begin{minipage}[b]{\linewidth}\sol{Salón y canales de comunicación de tareas.}\end{minipage} \\
			\hline
			\textbf{Situación} & \rule{0pt}{0.5cm}\begin{minipage}[b]{\linewidth}\sol{Circulan dos versiones de las indicaciones de una tarea.}\end{minipage} \\
			\hline
			\textbf{Necesidad / Problema / Incomodidad} & \rule{0pt}{0.5cm}\begin{minipage}[b]{\linewidth}\sol{Reconocer qué instrucciones están vigentes.}\end{minipage} \\
			\hline
			\textbf{¿A quién afecta?} & \rule{0pt}{0.45cm}\begin{minipage}[b]{\linewidth}\sol{Estudiantes y docentes.}\end{minipage} \\
			\hline
			\textbf{Evidencia} & \rule{0pt}{0.6cm}\begin{minipage}[b]{\linewidth}\sol{Escena simulada: se comparan dos mensajes distintos.}\end{minipage} \\
			\hline
			\textbf{¿Por qué importa?} & \rule{0pt}{0.5cm}\begin{minipage}[b]{\linewidth}\sol{Las dudas pueden ocasionar trabajo repetido.}\end{minipage} \\
			\hline
			\textbf{Posible solución} & \rule{0pt}{0.45cm}\begin{minipage}[b]{\linewidth}\sol{Abierta; requiere verificar si la dificultad se repite.}\end{minipage} \\
			\hline
		\end{tabular}
	\end{center}
	
	\vspace{0.15cm}
	
	\section*{Oportunidad 5}
	
	\begin{center}
		\small
		\begin{tabular}{|L{0.22\textwidth}|L{0.61\textwidth}|}
			\hline
			\textbf{Lugar} & \rule{0pt}{0.5cm}\begin{minipage}[b]{\linewidth}\sol{Patio, junto a recipientes para residuos.}\end{minipage} \\
			\hline
			\textbf{Situación} & \rule{0pt}{0.5cm}\begin{minipage}[b]{\linewidth}\sol{Hay dudas al clasificar un empaque después de comer.}\end{minipage} \\
			\hline
			\textbf{Necesidad / Problema / Incomodidad} & \rule{0pt}{0.5cm}\begin{minipage}[b]{\linewidth}\sol{Comprender dónde depositar cada residuo.}\end{minipage} \\
			\hline
			\textbf{¿A quién afecta?} & \rule{0pt}{0.45cm}\begin{minipage}[b]{\linewidth}\sol{Estudiantes y personal de aseo.}\end{minipage} \\
			\hline
			\textbf{Evidencia} & \rule{0pt}{0.6cm}\begin{minipage}[b]{\linewidth}\sol{Escena simulada: un estudiante cambia el empaque de recipiente.}\end{minipage} \\
			\hline
			\textbf{¿Por qué importa?} & \rule{0pt}{0.5cm}\begin{minipage}[b]{\linewidth}\sol{Puede afectar la separación de residuos.}\end{minipage} \\
			\hline
			\textbf{Posible solución} & \rule{0pt}{0.45cm}\begin{minipage}[b]{\linewidth}\sol{Abierta; queda como alternativa exploratoria.}\end{minipage} \\
			\hline
		\end{tabular}
	\end{center}
	
	\vspace{0.15cm}
	
	\section*{Oportunidad 6}
	
	\begin{center}
		\small
		\begin{tabular}{|L{0.22\textwidth}|L{0.61\textwidth}|}
			\hline
			\textbf{Lugar} & \rule{0pt}{0.5cm}\begin{minipage}[b]{\linewidth}\sol{Pasillo de acceso a salones durante un cambio de clase.}\end{minipage} \\
			\hline
			\textbf{Situación} & \rule{0pt}{0.5cm}\begin{minipage}[b]{\linewidth}\sol{Grupos detenidos coinciden con quienes necesitan pasar.}\end{minipage} \\
			\hline
			\textbf{Necesidad / Problema / Incomodidad} & \rule{0pt}{0.5cm}\begin{minipage}[b]{\linewidth}\sol{Circular sin bloquear el tránsito de otras personas.}\end{minipage} \\
			\hline
			\textbf{¿A quién afecta?} & \rule{0pt}{0.45cm}\begin{minipage}[b]{\linewidth}\sol{Estudiantes y docentes en desplazamiento.}\end{minipage} \\
			\hline
			\textbf{Evidencia} & \rule{0pt}{0.6cm}\begin{minipage}[b]{\linewidth}\sol{Escena simulada: una persona se detiene para poder pasar.}\end{minipage} \\
			\hline
			\textbf{¿Por qué importa?} & \rule{0pt}{0.5cm}\begin{minipage}[b]{\linewidth}\sol{Puede dificultar llegar con tranquilidad a clase.}\end{minipage} \\
			\hline
			\textbf{Posible solución} & \rule{0pt}{0.45cm}\begin{minipage}[b]{\linewidth}\sol{Abierta; no se investigará en esta secuencia.}\end{minipage} \\
			\hline
		\end{tabular}
	\end{center}
	
	\section*{Reflexión inicial}
	
	\textbf{1. ¿Qué oportunidad parece más interesante para investigar? ¿Por qué?}
	
	\vspace{0.1cm}
	
	\noindent\answerline{Elegimos investigar las filas para comprar alimentos: podrían limitar el descanso.}\rule{\linewidth}{0.4pt}
	
	\vspace{0.35cm}
	
	\noindent\answerline{La escena preliminar justifica observarlas; no conocemos aún sus causas.}\rule{\linewidth}{0.4pt}
	
	\vspace{0.4cm}
	
	\textbf{2. ¿Qué evidencia les hace pensar que esta es una situación real y no una suposición?}
	
	\vspace{0.1cm}
	
	\noindent\answerline{En el ejemplo simulado, la fila permanece mientras llegan más compradores.}\rule{\linewidth}{0.4pt}
	
	\vspace{0.35cm}
	
	\noindent\answerline{Es un indicio; todavía no demuestra una espera excesiva o evitable.}\rule{\linewidth}{0.4pt}
	
	\vspace{0.4cm}
	
	\textbf{3. ¿Quién podría ayudarles a comprender mejor esta situación?}
	
	\vspace{0.1cm}
	
	\noindent\answerline{Estudiantes que compran o desisten de hacerlo y personal del punto de venta.}\rule{\linewidth}{0.4pt}
	
	\vspace{0.35cm}
	
	\noindent\answerline{Sus experiencias permitirían contrastar lo observado, sin asumir una solución.}\rule{\linewidth}{0.4pt}
	
	\vspace{0.4cm}
	
	\textbf{4. ¿Qué necesitarían averiguar todavía?}
	
	\vspace{0.1cm}
	
	\noindent\answerline{Cuánto se espera, cómo cambia la fila y qué ocurre al elegir, pagar y recibir.}\rule{\linewidth}{0.4pt}
	
	\vspace{0.35cm}
	
	\noindent\answerline{Qué espera es evitable y qué parte responde a la demanda y capacidad del servicio.}\rule{\linewidth}{0.4pt}
	
	\section*{Verificación final}
	
	\noindent\checkbox\makebox[0pt][r]{\sol{\(\times\)}\hspace{0.15em}}\quad Identifiqué situaciones de mi entorno real.
	
	\vspace{0.15cm}
	
	\noindent\checkbox\makebox[0pt][r]{\sol{\(\times\)}\hspace{0.15em}}\quad Describí lo que realmente observé.
	
	\vspace{0.15cm}
	
	\noindent\checkbox\makebox[0pt][r]{\sol{\(\times\)}\hspace{0.15em}}\quad Identifiqué quiénes están afectados.
	
	\vspace{0.15cm}
	
	\noindent\checkbox\makebox[0pt][r]{\sol{\(\times\)}\hspace{0.15em}}\quad Incluí evidencia en lugar de basarme únicamente en opiniones.
	
	\vspace{0.15cm}
	
	\noindent\checkbox\makebox[0pt][r]{\sol{\(\times\)}\hspace{0.15em}}\quad Expliqué por qué la situación podría ser importante.
	
	\vspace{0.15cm}
	
	\noindent\checkbox\makebox[0pt][r]{\sol{\(\times\)}\hspace{0.15em}}\quad No asumí una solución final demasiado pronto.
	
	\vspace{0.3cm}
	
	\begin{center}
		\textbf{Primero comprendan la situación. Luego comprendan a las personas.}
	\end{center}

	
\end{document}
